\chapter*{แผนบริหารการสอนประจำวิชา}
\addcontentsline{toc}{chapter}{แผนบริหารการสอนประจำวิชา}

\begin{tcolorbox}[enhanced, colback=white, colframe=waterNavy, colbacktitle=waterNavy,
    coltitle=white, fonttitle=\bfseries\normalsize,
    title={ข้อมูลทั่วไปของรายวิชา},
    arc=2mm, boxrule=0.8pt, leftrule=4mm,
    top=8pt, bottom=8pt, left=10pt, right=10pt]
\begin{tabularx}{\textwidth}{l>{\raggedright\arraybackslash}X}
\textbf{รหัสวิชา} & MIC3205 \\
\textbf{ชื่อรายวิชา} & จุลชีววิทยาสิ่งแวดล้อมและการวิเคราะห์คุณภาพน้ำ \\
 & (Environmental Microbiology and Water Quality Analysis) \\
\textbf{จำนวนหน่วยกิต} & 3(2-2-5) หน่วยกิต (บรรยาย 2 ชั่วโมง ปฏิบัติการ 2 ชั่วโมง ศึกษาด้วยตนเอง 5 ชั่วโมง) \\
\textbf{หลักสูตร} & วิทยาศาสตรบัณฑิต สาขาวิชาจุลชีววิทยา / ดิจิทัลการเกษตรและสิ่งแวดล้อม \\
\textbf{อาจารย์ผู้สอน} & ผู้ช่วยศาสตราจารย์ ดร.จิรภัทร จันทมาลี และ ผู้ช่วยศาสตราจารย์ ดร.ชีวะ ทัศนา \\
\textbf{ภาคการศึกษา/ชั้นปี} & ภาคการศึกษาที่ 1 ชั้นปีที่ 3 \\
\end{tabularx}
\end{tcolorbox}

\begin{tcolorbox}[enhanced, colback=cleanSky!25!white, colframe=oceanTeal, colbacktitle=oceanTeal,
    coltitle=white, fonttitle=\bfseries\normalsize,
    title={คำอธิบายรายวิชา (Course Description)},
    arc=2mm, boxrule=0.8pt, leftrule=4mm,
    top=8pt, bottom=8pt, left=10pt, right=10pt, before skip=8pt]
ศึกษาหลักความปลอดภัยและมาตรฐานการปฏิบัติการในห้องปฏิบัติการทางชีวภาพ การแพร่กระจายและบทบาทของจุลินทรีย์ในระบบนิเวศ ดิน น้ำ และบรรยากาศ การสำรวจประชากรจุลินทรีย์ในแหล่งน้ำธรรมชาติและน้ำเสีย เทคนิคการย้อมสีและสัณฐานวิทยาของเซลล์ เทคนิคการแยกเชื้อบริสุทธิ์และการเพาะเลี้ยง การเตรียมอาหารเลี้ยงเชื้อ จลนศาสตร์การเจริญเติบโตของจุลินทรีย์ จุลินทรีย์กลุ่มไร้อากาศในระบบบำบัดน้ำ เมแทบอลิซึมและการทดสอบชีวเคมี การควบคุมและการยับยั้งจุลินทรีย์ และการประยุกต์ใช้นวัตกรรมไบโอบล็อกและจุลินทรีย์ตรึงรูปในการฟื้นฟูคุณภาพน้ำทางชีวภาพ
\end{tcolorbox}

\vspace{0.4cm}

\section*{แผนการสอนรายสัปดาห์ (15 สัปดาห์)}

\noindent\begin{tabularx}{\textwidth}{|c|>{\raggedright\arraybackslash}X|c|}
\hline
\rowcolor{waterNavy}
\textcolor{white}{\textbf{สัปดาห์}} & \textcolor{white}{\textbf{หัวข้อบรรยายและหัวข้อปฏิบัติการ}} & \textcolor{white}{\textbf{ชั่วโมง}} \\
\hline
1 & ความปลอดภัยในห้องปฏิบัติการและเครื่องมือพื้นฐานทางจุลชีววิทยา & 4 \\
\hline
2 & การแพร่กระจายของจุลินทรีย์ในสิ่งแวดล้อมและบรรยากาศ & 4 \\
\hline
3 & การเก็บตัวอย่างและการสำรวจจุลินทรีย์ในแหล่งน้ำธรรมชาติ & 4 \\
\hline
4 & การใช้กล้องจุลทรรศน์และเทคนิคการย้อมสีแบคทีเรีย (Simple \& Gram Stain) & 4 \\
\hline
5 & สัณฐานวิทยาและโครงสร้างพิเศษของเซลล์จุลินทรีย์ (Endospore \& Capsule) & 4 \\
\hline
6 & เทคนิคการแยกเชื้อบริสุทธิ์ (Streak Plate, Pour Plate, Spread Plate) & 4 \\
\hline
7 & อาหารเลี้ยงเชื้อ การจำแนกประเภท และการเตรียมอาหารสังเคราะห์ & 4 \\
\hline
\rowcolor{cleanSky!50!white}
\textbf{8} & \textbf{สอบประเมินผลกลางภาค (Midterm Examination)} & \textbf{2} \\
\hline
9 & จลนศาสตร์การเจริญเติบโตของจุลินทรีย์และการวัดชีวมวล & 4 \\
\hline
10 & จุลินทรีย์ไม่ใช้ออกซิเจนและการเพาะเลี้ยงในสภาวะแอนแอโรบิก & 4 \\
\hline
11 & เมแทบอลิซึมและการทดสอบคุณสมบัติทางชีวเคมี (IMViC Test, Catalase) & 4 \\
\hline
12 & การควบคุมจุลินทรีย์ด้วยวิธีทางกายภาพและสารเคมีฆ่าเชื้อ & 4 \\
\hline
13 & แบคทีเรียกลุ่มโคลิฟอร์มและการตรวจวัดดัชนีชี้วัดคุณภาพน้ำทางจุลชีววิทยา (MPN) & 4 \\
\hline
14 & นวัตกรรมไบโอบล็อก (Bioblock) และระบบบำบัดน้ำเสียทางชีวภาพ & 4 \\
\hline
15 & การนำเสนอโครงงานวิจัยการตรวจวิเคราะห์คุณภาพน้ำภาคสนาม & 4 \\
\hline
\rowcolor{cleanSky!50!white}
\textbf{16} & \textbf{สอบประเมินผลปลายภาค (Final Examination)} & \textbf{2} \\
\hline
\end{tabularx}

\clearpage
